% =====================================================================
%  Hiện thực UC53, UC48, UC49 (chèn vào hien-thuc-uc.tex trước mục
%  "Các use case còn lại")
% =====================================================================
\subsection{Cấu hình hệ thống và khóa tài khoản (UC53, UC48)}\label{sssec:ht-quan-tri}

Hai thao tác quản trị có tác động toàn hệ thống thuộc module admin: \uc{UC53} qua D09 \texttt{PUT} \path{/admin/configs} (\texttt{admin.Service.UpdateConfigs}) và thao tác khóa tài khoản của \uc{UC48} qua D02 \texttt{PATCH} \path{/admin/users/{userId}} (\texttt{admin.Service.DisableUser}). Thay đổi của cả hai phải lan tới mọi tiến trình đang chạy chứ không dừng ở cơ sở dữ liệu: tham số mới có hiệu lực ở lần xử lý kế tiếp trên mọi bản API và worker, tài khoản bị khóa mất quyền truy cập ngay. Hình~\ref{fig:seq-quan-tri-cau-hinh} thể hiện hai luồng này.

\diagram[\textwidth]{c3-38-seq-quan-tri-cau-hinh}{Biểu đồ tuần tự hiện thực use case Cấu hình tham số hệ thống và thao tác khóa tài khoản}{fig:seq-quan-tri-cau-hinh}

Bộ định tuyến chặn nhóm đường dẫn \path{/admin/*} với người gọi có \texttt{system\_role} khác \texttt{admin} (403 mã 3007, \msg{MSG13}). D09 nhận ánh xạ từ khóa tham số sang giá trị và ghi tất cả trong một giao dịch, vì \uc{UC53} cấm áp dụng một phần và ba trọng số gợi ý chỉ có nghĩa khi đổi cùng nhau. Các dòng \texttt{system\_configs} liên quan được khóa theo thứ tự khóa (bước 4), nên hai quản trị viên lưu đồng thời được tuần tự hóa mà không khóa chết. Trình kiểm tra là mã chung cho mọi khóa vì kiểu và miền hợp lệ nằm sẵn trong các cột \texttt{value\_type}, \texttt{min\_value}, \texttt{max\_value} (mục~\ref{sec:csdl}); riêng \texttt{suggest.weights} được kiểm thêm tổng bằng 1 bằng số thập phân chính xác, để $0{,}6 + 0{,}3 + 0{,}1$ không bị số thực dấu phẩy động đánh giá sai (bước 5). Một khóa lạ hay một giá trị sai hủy cả giao dịch với 400 mã 2002 kèm \msg{MSG27}, mảng \texttt{errors} nêu từng khóa vi phạm cùng miền hợp lệ (bước 6--7). Khi hợp lệ, mọi khóa được cập nhật cùng đúng một dòng \texttt{audit\_logs} chứa giá trị cũ và mới (bước 8). D09 không dùng \texttt{If-Match} vì cấu hình không có cột phiên bản: hai lần lưu đồng thời thì lần sau thắng, nhưng nhật ký D10 cho thấy và hoàn tác được.

Mỗi tiến trình giữ cấu hình trong bộ nhớ, nạp từ khóa Redis \texttt{cfg:all} sống 5 phút và chỉ đọc PostgreSQL khi khóa này trống. Sau commit, dịch vụ xóa \texttt{cfg:all}, xóa \texttt{suggest:*} nếu tham số \texttt{suggest.*} đổi, rồi phát \texttt{config.changed} trên kênh \texttt{cfg} (bước 9); mọi bản API và worker bỏ bản trong bộ nhớ và nạp lại ở lần đọc kế tiếp (bước 10). Pub/Sub giao tối đa một lần, nên thời hạn sống 5 phút là cận trên của độ cũ khi thông điệp bị mất. “Lần xử lý kế tiếp” của \uc{UC53} vì thế có nghĩa cụ thể: \texttt{ReminderDispatcher} đọc lại \texttt{reminder.runAt} ở lần kiểm tra mỗi phút, nếu hôm nay đã chạy thì khung giờ mới áp dụng từ hôm sau; \texttt{otp.ttlMinutes}, \texttt{invite.ttlHours} chỉ áp dụng cho mã sinh sau thay đổi; giảm \texttt{group.maxMembers} hay \texttt{mealplan.horizonDays} không loại thành viên hay món đã có.

D02 gồm khóa, mở khóa, cấp và thu hồi quyền quản trị; khóa có hệ quả lớn nhất nên được vẽ trong hình. Module admin không sở hữu bảng tài khoản nên gọi cổng của identity trong cùng giao dịch: khóa dòng \texttt{users}, từ chối tự khóa chính mình (422 mã 4004), đặt \texttt{status = 'disabled'} cùng \texttt{lock\_reason}, tăng \texttt{token\_version}, thu hồi mọi phiên, ghi \texttt{audit\_logs} và chèn \texttt{SendEmail} báo chủ tài khoản (bước 14--17). Sau commit, \path{auth:ver:<userId>} được ghi đè bằng phiên bản mới thay vì chỉ xóa, nên mọi bản API từ chối ngay mã truy cập cũ bằng mã 3002 mà không phải đọc CSDL; lần làm mới A05 kế tiếp nhận 403 mã 3005 và ứng dụng hiển thị \msg{MSG21}. Kênh \texttt{user:\{uid\}} báo bản API đang giữ WebSocket của người bị khóa đóng kết nối (bước 18--19). Mở khóa đưa tài khoản về \texttt{active}, xóa \texttt{lock\_reason}, không khôi phục phiên cũ. Cấp hay thu hồi quyền quản trị cũng tăng \texttt{token\_version}, vì vai trò nằm trong claim \texttt{role} của JWT. Thu hồi quyền khóa mọi dòng quản trị viên đang hoạt động theo \texttt{id} rồi mới đếm, nên hai quản trị viên cuối cùng không thể đồng thời thu hồi quyền của nhau (\uc{UC48} bước 3b). D01 chỉ trả email, họ tên, trạng thái, ngày tạo và \texttt{last\_login\_at}; quản trị viên không là thành viên nên không đọc được dữ liệu nhóm (\br{BR23}).

% =====================================================================
\subsection{Duyệt thực phẩm do người dùng tạo (UC49)}\label{sssec:ht-duyet-thuc-pham}

Thực phẩm riêng tạo ở \uc{UC25} bước 3a và \uc{UC26} bước 4a mang \texttt{scope = space}, \texttt{status = pending\_review} và chỉ xuất hiện trong không gian đã tạo ra nó. Phần duyệt của \uc{UC49} khép lại vòng đời đó: quản trị viên chuẩn hóa thực phẩm thành bản dùng chung, hoặc gộp nó vào một bản dùng chung có sẵn. \texttt{catalog.Service.ReviewFood} hiện thực việc này qua D11 \texttt{POST} \path{/admin/foods/{foodId}/review} với thân \texttt{\{action, targetFoodId\}}, \texttt{action} là \texttt{approve} hoặc \texttt{merge}; hàng chờ lấy từ D03 lọc \texttt{status = pending\_review}. Hình~\ref{fig:seq-duyet-thuc-pham} thể hiện hai nhánh cùng phần xử lý nền sau commit.

\diagram[\textwidth]{c3-39-seq-duyet-thuc-pham}{Biểu đồ tuần tự hiện thực thao tác duyệt và gộp thực phẩm riêng của use case Quản lý danh mục thực phẩm}{fig:seq-duyet-thuc-pham}

D11 dùng \texttt{Idempotency-Key} vì quản trị viên cũng thao tác trên mạng di động. Giao dịch khóa thực phẩm nguồn và đích theo \texttt{id} tăng dần (bước 3), nên hai quản trị viên duyệt cùng một thực phẩm hay gộp chéo hai thực phẩm vào nhau không khóa chết; người đến sau thấy thực phẩm không còn \texttt{pending\_review} và nhận 409 mã 4002 kèm \msg{MSG35}. Nhánh \texttt{approve} đổi phạm vi thành \texttt{global}, kích hoạt và bỏ \texttt{owner\_space\_id} (bước 6). Hàng rào là hai chỉ mục duy nhất theo phạm vi: tên trong cùng loại (\br{BR22}) và mã vạch là duy nhất trong phạm vi \texttt{global}, nên thực phẩm riêng trùng tên hoặc mã vạch với một bản dùng chung không thể thành bản dùng chung thứ hai. Dịch vụ kiểm tra trước để trả 409 mã 4003 kèm \msg{MSG02} và gợi ý gộp (bước 5); lỗi 23505 khi có truy cập đồng thời được ánh xạ về cùng mã. Nhánh \texttt{merge} yêu cầu đích là bản dùng chung đang hoạt động, khác nguồn, ngược lại trả 400 mã 2002 (bước 7); nguồn chuyển \texttt{inactive} với \texttt{merged\_into} trỏ tới đích (bước 8). Cả hai nhánh ghi \texttt{audit\_logs} và chèn \texttt{IndexFood} cho \texttt{food.changed} trong cùng giao dịch (bước 9).

Gộp không sửa trực tiếp dữ liệu tham chiếu, vì pantry, shopping và recipe phụ thuộc catalog chứ không có chiều ngược lại; catalog gọi vào ba module đó sẽ tạo chu trình (mục~\ref{sec:kien-truc}). Sự kiện \texttt{food.merged} vì vậy sinh tác vụ \texttt{RemapFood} cho từng module nhận (bước 11). Thực phẩm riêng chỉ được tham chiếu trong không gian sở hữu, nên mỗi tác vụ chỉ chạm vài chục đến vài trăm dòng của một không gian và chạy trong một giao dịch: đổi \texttt{food\_id} của lô, mặt hàng hay nguyên liệu sang đích, tăng \texttt{version}, ghi \texttt{space\_changes}, phát \texttt{space:\{sid\}} sau commit. Tác vụ lũy đẳng vì câu \texttt{UPDATE ... WHERE food\_id = <nguồn>} chạy lại không còn dòng nào để đổi. Riêng nguyên liệu, khóa chính (\texttt{recipe\_id}, \texttt{food\_id}) không cho công thức đã có thực phẩm đích nhận dòng thứ hai, nên lượng của dòng nguồn được cộng vào dòng đích sau khi quy đổi nếu cùng nhóm đại lượng (\br{BR16}), ngược lại dòng nguồn được giữ để người dùng tự xử lý; hai mặt hàng trùng thực phẩm trong một danh sách không bị tự gộp theo \br{BR09}. Trong lúc tác vụ chưa xong, cột \texttt{merged\_into} cho phép máy chủ chuyển lệnh ngoại tuyến còn mang id cũ sang thực phẩm đích thay vì trả 404. \texttt{IndexFood} xóa tài liệu Elasticsearch của nguồn khi gộp, hoặc ghi lại tài liệu với \texttt{scope = global} khi duyệt để thực phẩm xuất hiện trong gợi ý của mọi không gian (bước 12, Hình~\ref{fig:seq-tim-kiem-chi-muc}).
